% 原书第 206—209 页；第 13.4—13.5 节全文。
\section{暴露机制的协同估计}
\label{sec:ch13-collaborative-exposure}

尽管干预 $A$ 采用平衡分配的随机化，估计已知的暴露机制 $g_0(A\mid W)=0.5$，仍可提高分析精度（van der Laan and Robins，2003）。与前面一样，我们既希望遵循研究设计，即是否采用配对匹配，也希望只在纳入某个协变量能够提高经验效率时，才对它作调整。例如，对于在干预组之间不平衡、即能够预测 $A$，却不能预测结局的协变量，一般不希望对其作调整。同样，如果某个协变量，如 $W1$，已经纳入条件结局均值 $\bar Q_0(A,W)$ 的工作模型，那么在估计暴露机制时进一步调整它，可能并不能提高精度。为此，我们在算法中纳入 C-TMLE 方法，见第 10 章。

\noindent\textbf{步骤 2 的预先指定自适应方法：靶向更新。}
首先，提出暴露机制 $g_0(A\mid W)$ 的候选估计量库。如前所述，应预先在研究方案或分析计划中指定这一算法库。一个可能的算法库可以由以下 logistic 回归工作模型组成：
\[
 \begin{aligned}
 \logit[g^{(a)}(W)]&=\beta_0,\\
 \logit[g^{(b)}(W)]&=\beta_0+\beta_1W1,\\
 \logit[g^{(c)}(W)]&=\beta_0+\beta_1W2,
 \end{aligned}
\]
例如，$W1$ 为基线流行率，$W2$ 为男性包皮环切覆盖率。对于给定的条件结局均值初始估计量 $\bar Q_n(A,W)$，每个算法都会产生不同的更新；这个初始估计量，由前述适用于无匹配或有匹配试验的步骤 1 数据自适应方法选出。换言之，$g_0(A\mid W)$ 的每个候选估计量，都产生一个不同的靶向估计量 $\bar Q_n^*(A,W)$。还应注意，第一个工作模型对应于未经调整的估计量。

为了在候选算法之间选择，需要预先指定一个表现度量。如前所述，我们建议，以 TMLE 的估计渐近方差作为风险，并使其与研究设计，即是否采用配对匹配，以及科学问题，即总体或样本效应相适应。因此，所采用的损失函数为：
\begin{itemize}
 \item 无匹配、以 PATE 为目标：式 (13.5) 中的 $\mathcal L^{\mathcal P}(g,\bar Q_n)$；
 \item 无匹配、以 SATE 为目标：式 (13.6) 中的 $\mathcal L^S(g,\bar Q_n)$；
 \item 有匹配、以 PATE 为目标：式 (13.8) 中的 $\bar{\mathcal L}^{\mathcal P}(g,\bar Q_n)$；
 \item 有匹配、以 SATE 为目标：式 (13.9) 中的 $\bar{\mathcal L}^S(g,\bar Q_n)$。
\end{itemize}
其中，$g$ 表示暴露机制的一个候选估计量，$\bar Q_n$ 表示已经选出的结局回归初始估计量。

最后，需要预先指定交叉验证方案，以得到真实反映风险的度量，并减小过拟合的可能性。与前面一样，我们介绍 $V$ 折交叉验证，将数据划分为 $V$ 折，每折大小约为 $n/V$。如果采用了匹配，划分时应保留匹配对。用于协同估计暴露机制的交叉验证选择器，可以按如下步骤实施。
\begin{enumerate}[label=(\Alph*)]
 \item 依次对每一折 $v=\{1,\ldots,V\}$：
 \begin{enumerate}[label=(\alph*)]
  \item 将第 $v$ 折中的观测作为验证集，其余观测作为训练集。
  \item 仅使用训练集数据，拟合结局回归的初始估计量 $\bar Q_n(A,W)$。
  \item 仅使用训练集数据，拟合估计暴露机制的每个算法。对于上述算法库，根据工作模型，将暴露 $A$ 对协变量 $W$ 作 logistic 回归。将估计的暴露机制分别记为
  $g_n^{(a)}(A\mid W)$、$g_n^{(b)}(A\mid W)$ 和 $g_n^{(c)}(A\mid W)$。
  \item 对每个算法，使用其估计的暴露机制拟合，对初始估计量 $\bar Q_n(A,W)$ 进行靶向更新。将靶向回归拟合分别记为
  $\bar Q_n^{(a),*}(A,W)$、$\bar Q_n^{(b),*}(A,W)$ 和 $\bar Q_n^{(c),*}(A,W)$，其中，上标对应于估计暴露机制所用的算法。
  \item 对每个算法，为验证集观测分别在处理和对照条件下的结局，得到靶向预测。例如，对于拟合暴露机制的第一个算法，验证集观测 $O_k$ 的预测为
  $\bar Q_n^{(a),*}(1,W_k)$ 和 $\bar Q_n^{(a),*}(0,W_k)$。
  \item 对每个算法，代入其特定的预测，评价验证集观测的损失函数。例如，如果推断目标是无匹配试验中的 SATE，第一个算法对于验证集观测 $O_k$ 给出
  \[
   \begin{aligned}
   &\mathcal L^S(g_n^{(a)},\bar Q_n^{(a),*})(O_k)\\
   &\quad=\left[\left\{\frac{\mathbb I(A_k=1)}{g_n^{(a)}(1\mid W_k)}
              -\frac{\mathbb I(A_k=0)}{g_n^{(a)}(0\mid W_k)}\right\}
         (Y_k-\bar Q_n^{(a),*}(A_k,W_k))\right]^2.
   \end{aligned}
  \]
  \item 对每个算法，将验证集 $v$ 中各观测的估计损失取平均，得到风险估计。例如，如果推断目标是无匹配试验中的 SATE，那么对于估计暴露机制的第一个算法，有
  \[
   \mathrm{Risk}_v^{(a)}
   =\frac1{n_v}\sum_{k\in v}
     \mathcal L^S(g_n^{(a)},\bar Q_n^{(a),*})(O_k),
  \]
  其中，$n_v$ 为验证集 $v$ 中的观测数。
 \end{enumerate}
 \item 对每个算法，将 $V$ 折的风险估计取平均。
 \item 选择交叉验证风险最小的算法。这个算法给出最小的交叉验证方差估计。
\end{enumerate}
随后，用所选估计量实施 TMLE 算法步骤 2 中的靶向更新。

在这个方案中，我们将结局回归的初始估计量 $\bar Q_n(A,W)$ 视为固定，提出第二轮交叉验证，选择暴露机制的拟合。另一种做法是，构造一个候选 TMLE 库，以结局回归初始估计量和暴露机制估计量的选择为索引，再以影响曲线的交叉验证方差作为表现度量，从该算法库中选择。在整群随机试验中，我们推荐双重交叉验证方法（double cross-validation），以遵循估计暴露机制的协同原则，并避免过拟合。只有在估计 $\bar Q_0(A,W)$ 时所作的调整之外，还能进一步提高效率时，我们才希望估计 $g_0(A\mid W)$。双重交叉验证方法根据结局回归 $\bar Q_0(A,W)$ 的拟合，来估计暴露机制 $g_0(A\mid W)$。例如，可以根据对 $\bar Q_0(A,W)$ 的选择，限制 $g_0(A\mid W)$ 的算法库：如果已选择某个协变量用于结局回归估计，就将对应算法从暴露机制的算法库中移除。最后，双重交叉验证方法允许考虑大量可能的候选 TMLE，即 $\bar Q_0(A,W)$ 估计量与 $g_0(A\mid W)$ 估计量的所有可能组合。在拥有许多条件独立单位的试验中，可以考虑单轮交叉验证方法，它对应于离散超级学习器；也可以采用完整超级学习器，以影响曲线平方作为损失函数。

\section{统计推断}
\label{sec:ch13-inference}

概括而言，我们提出以下数据自适应 C-TMLE，以使随机试验的精度与功效最大化。
\begin{itemize}
 \item 步骤 1：用工作模型 $\bar Q_n(A,W)$ 初始估计条件结局均值。该模型通过数据自适应选择，使无匹配试验或有匹配试验的分析经验效率最大化。
 \item 步骤 2：使用估计的暴露机制 $g_n(A\mid W)$，对初始估计量 $\bar Q_n(A,W)$ 进行靶向更新。暴露机制也是通过数据自适应选择，以进一步使分析的经验效率最大化。
 \item 步骤 3：对靶向预测之差取平均，得到点估计：
 \[
  \Psi_n(\bar Q_n^*)
  =\frac1n\sum_{i=1}^n[\bar Q_n^*(1,W_i)-\bar Q_n^*(0,W_i)].
 \]
\end{itemize}

现在，需要一个考虑选择过程的方差估计量。为此，我们提出使用交叉验证方差估计量。与前面一样，遵循条件独立单位，将数据分为验证集和训练集。使用训练集数据拟合所选 TMLE，并用它估计验证集观测的影响曲线\footnote{对于配对匹配试验中总体效应的 TMLE，还需要修正项 $\rho_n$ 的交叉验证估计，见式 (13.7)。这个项是残差的函数；利用训练集拟合靶向估计量 $\bar Q_n^*(A,W)$，就可以为验证集中的每个匹配对估计该项。}。随后，可用影响曲线交叉验证估计的样本方差，进行假设检验并构造 Wald 型置信区间。应当注意，为了评价候选估计量的表现，我们已经计算了影响曲线的交叉验证估计。因此，这一步不需要额外计算；选择过程已经给出了交叉验证方差的一个估计。还应注意，对于较小的算法库，例如只有两个候选 TMLE，模拟结果支持采用标准方差估计量，而非交叉验证方差估计量来进行推断（Balzer et al.，2016b）。

